\documentclass{article}
\newtheorem{theorem}{Theorem}[section]
\newtheorem{corollary}[theorem]{Corollary}

\title{A vanishing theorem for weighted stable curves}
\author{F. Ekstrom \and H. Gallant}

\begin{document}
\maketitle

\section{The vanishing}

\begin{theorem}\label{thm:vanish}
Let $C$ be a weighted stable curve of genus $g$ with weights summing to less than one. Then $H^1(C, \omega_C) = 0$.
\end{theorem}

\begin{proof}
The weight condition forces every component to be rational, so the dualising sheaf is a sum of line bundles of negative degree. The vanishing for each summand is the local computation of \cite[Theorem 1.1]{other}, and the compatibility of the two normalisations is the point of \cite{third}.
\end{proof}

\section{A consequence}

\begin{corollary}\label{cor:rigid}
Such a curve admits no infinitesimal automorphism.
\end{corollary}

\begin{proof}
Immediate from Theorem~\ref{thm:vanish} and Serre duality.
\end{proof}

\begin{thebibliography}{9}
\bibitem{other} B. Aurelio and D. Cerny, \emph{Tropical cycles on fans of valuated matroids}, preprint, 2024, arXiv:2401.00001.
\bibitem{third} J. Ibarra, \emph{Foundations of the synthetic period map}, Lecture notes, 1998.
\end{thebibliography}

\end{document}
